% ============================================================
\chapter{Forecasting and Confidence Intervals}
\label{ch:prevision}
% ============================================================

\section{Motivating Example}
Forecasting electricity demand for the next day, together with a confidence interval, is essential for grid management. The quality of the forecast determines economic costs.

\section{Optimal Forecast}\index{forecasting}

\begin{definition}[Optimal Forecast]
The \textbf{optimal forecast} of $X_{n+h}$ given $X_1,\ldots,X_n$ in the least squares sense is:
\[
\hat{X}_{n+h|n} = \E[X_{n+h}|X_1,\ldots,X_n].
\]
The forecast error is $e_{n+h|n}=X_{n+h}-\hat{X}_{n+h|n}$.
\end{definition}

\begin{theorem}[Mean Squared Error]
The conditional expectation minimizes the MSE:
\[
\hat{X}_{n+h|n} = \arg\min_{g(X_1,\ldots,X_n)}\E\left[(X_{n+h}-g)^2\right].
\]
\end{theorem}

\section{ARMA Forecasting}

\begin{theorem}[Forecast of an ARMA$(p,q)$]
Using the MA$(\infty)$ representation: $X_t=\sum_{j=0}^{\infty}\psi_j\eps_{t-j}$, the $h$-step-ahead forecast is:
\[
\hat{X}_{n+h|n} = \sum_{j=h}^{\infty}\psi_j\eps_{n+h-j}.
\]
The forecast error is $e_{n+h|n}=\sum_{j=0}^{h-1}\psi_j\eps_{n+h-j}$, with variance:
\[
\sigma_h^2 = \sigma^2\sum_{j=0}^{h-1}\psi_j^2.
\]
\end{theorem}

\begin{corollary}
For a stationary ARMA, $\sigma_h^2\to\gamma(0)=\Var(X_t)$ as $h\to\infty$: the forecast converges to the unconditional mean and the confidence interval widens to the unconditional interval.
\end{corollary}

\section{Prediction Intervals}\index{prediction interval}

\begin{definition}
Under the Gaussian assumption, a $(1-\alpha)$ prediction interval for $X_{n+h}$ is:
\[
\hat{X}_{n+h|n} \pm z_{1-\alpha/2}\,\sigma_h,
\]
where $z_{1-\alpha/2}$ is the quantile of the standard normal distribution.
\end{definition}

\begin{warning}
Prediction intervals do not account for the uncertainty in estimated parameters. Bootstrap or simulation methods can be used to obtain more realistic intervals.
\end{warning}

\section{ARIMA Forecasting}

\begin{theorem}
For an ARIMA$(p,d,q)$, the long-term forecast of the \emph{level} $X_{n+h}$ grows linearly (if $d=1$) or quadratically (if $d=2$): the prediction interval widens without bound.
\end{theorem}

\section{Evaluation Criteria}\index{evaluation}

\begin{definition}[Forecast Metrics]
\begin{align}
\text{MAE} &= \frac{1}{H}\sum_{h=1}^H|e_{n+h|n}|, \\
\text{RMSE} &= \sqrt{\frac{1}{H}\sum_{h=1}^H e_{n+h|n}^2}, \\
\text{MAPE} &= \frac{100}{H}\sum_{h=1}^H\frac{|e_{n+h|n}|}{|X_{n+h}|}\%.
\end{align}
\end{definition}

\begin{definition}[Temporal Cross-Validation]\index{cross-validation}
To evaluate predictive ability, one uses a \textbf{train/test} split that respects the temporal order (no shuffling). The \textbf{rolling window} or \textbf{expanding window} method is standard.
\end{definition}

\section{Diebold--Mariano Test}\index{test!Diebold--Mariano}

\begin{definition}
To compare two forecasting models, let $d_t = L(e_{1,t})-L(e_{2,t})$ be the loss differential. Under $H_0$ (equality):
\[
\text{DM} = \frac{\bar{d}}{\sqrt{\hat\sigma_d^2/H}} \dto \mathcal{N}(0,1).
\]
\end{definition}

\section{Implementation}

\begin{codeblock}[title=\textbf{Python}]
\begin{minted}{python}
import numpy as np
import pandas as pd
import matplotlib.pyplot as plt
from statsmodels.tsa.arima.model import ARIMA
from sklearn.metrics import mean_absolute_error, mean_squared_error

# Simuler un ARMA(1,1)
np.random.seed(42)
from statsmodels.tsa.arima_process import ArmaProcess
ar = np.array([1, -0.7])
ma = np.array([1, 0.3])
data = ArmaProcess(ar, ma).generate_sample(nsample=300)

# Train/test split
train, test = data[:250], data[250:]

# Ajuster et prévoir
model = ARIMA(train, order=(1, 0, 1))
results = model.fit()

forecast = results.get_forecast(steps=50)
pred = forecast.predicted_mean
ci = forecast.conf_int(alpha=0.05)

# Métriques
mae = mean_absolute_error(test, pred)
rmse = np.sqrt(mean_squared_error(test, pred))
print(f"MAE  = {mae:.4f}")
print(f"RMSE = {rmse:.4f}")

# Graphique
fig, ax = plt.subplots(figsize=(12, 5))
ax.plot(range(250), train, 'b-', lw=0.7, label='Train')
ax.plot(range(250, 300), test, 'g-', lw=1, label='Test')
ax.plot(range(250, 300), pred, 'r--', lw=1.5, label='Prévision')
ax.fill_between(range(250, 300), ci[:, 0], ci[:, 1],
                alpha=0.2, color='red')
ax.legend()
ax.set_title('Prévision ARMA(1,1)')
plt.savefig('ch10_forecast.pdf')
plt.show()

# Rolling forecast
rolling_preds = []
for i in range(len(test)):
    train_i = data[:250+i]
    model_i = ARIMA(train_i, order=(1,0,1)).fit()
    rolling_preds.append(model_i.forecast(steps=1)[0])
rmse_roll = np.sqrt(mean_squared_error(test, rolling_preds))
print(f"RMSE rolling: {rmse_roll:.4f}")
\end{minted}
\end{codeblock}

\section{Exercises}

\begin{exercise}[$\star$]
For an AR(1) with $\phi=0.8$ and $\sigma^2=1$, compute $\hat{X}_{n+1|n}$, $\hat{X}_{n+2|n}$ and the corresponding forecast error variances.
\end{exercise}

\begin{exercise}[$\star\star$]
Show that the $h$-step-ahead forecast error variance of an AR(1) is $\sigma^2\frac{1-\phi^{2h}}{1-\phi^2}$.
\end{exercise}

\begin{exercise}[$\star\star$]
Explain why the MAPE is problematic when $X_t$ takes values close to zero.
\end{exercise}

\begin{exercise}[$\star\star\star$ -- Project]
Compare out-of-sample forecasts of an ARIMA, exponential smoothing, and a naive model ($\hat{X}_{n+h}=X_n$) on 5 real-world series. Use the Diebold--Mariano test.
\end{exercise}

\begin{keyformulas}
\begin{align*}
\hat{X}_{n+h|n} &= \sum_{j=h}^{\infty}\psi_j\eps_{n+h-j}, \quad \sigma_h^2=\sigma^2\sum_{j=0}^{h-1}\psi_j^2 \\
\text{CI}_{95\%} &: \hat{X}_{n+h|n}\pm 1.96\,\sigma_h \\
\text{RMSE} &= \sqrt{\frac{1}{H}\sum e_{n+h}^2}
\end{align*}
\end{keyformulas}
